\documentclass{article}
\usepackage{graphicx} % Required for inserting images
\usepackage{amsmath,amssymb,amsfonts}
\usepackage{xcolor}
\usepackage{mathtools}
\usepackage{hyperref}
\usepackage{blindtext}




\title{Quadratic approximation of determinant of deformation gradient in MPM}
\date{\today}
\begin{document}
\newcommand{\dd}{\partial}
\renewcommand{\vec}[1]{\mathbf{#1}}


\maketitle
We know from this  \href{https://www.cs.ucr.edu/~craigs/papers/2016-mpm-course/mpmcourse.pdf}{MPM course}\href and basically any MPM paper that:
\[F^{n+1}_{p} = (I+\Delta t\nabla \vec{v}) F_p^{n}\]
if we define $J^{n+1}=\det(F_p^{n+1})$ then $J^{n+1} = J^n\det(I+\Delta t\nabla \vec{v})$. \\
\textbf{Note:} From taylor expansion we know $e^x = 1 + x + \frac{x^2}{2!}+\frac{x^3}{3!}+\cdots$\\
Same is true for matrices, we can write $e^A = I + A + \frac{A^2}{2!}+\frac{A^3}{3!}+\cdots$\\

So we can write $I+\Delta t\nabla \vec{v} \approx e^{\Delta t\nabla \vec{v}}$\\
\textbf{Jacobi formula} tells us $\det(e^A) = e^{\text{tr}(A)}$, which can be further approximated using taylor approximation as $e^{\text{tr}(A)} \approx 1 + \text{tr}(A) + \frac{\text{tr}(A)^2}{2!}+\cdots$\\
so we can write: 
\begin{align*}
J^{n+1} &= J^n\det(I+\Delta t\nabla \vec{v})\\
&\approx J^n\det(e^{\Delta t\nabla \vec{v}})\\
&= J^n e^{\text{tr}(\Delta t\nabla \vec{v})}\\
&\approx J^n\left(1 + \text{tr}(\Delta t\nabla \vec{v}) + \frac{\text{tr}(\Delta t\nabla \vec{v})^2}{2!}\right)
\end{align*}
in normal circumstances, we would just use $J^{n+1} = J^n\det(I+\Delta t\nabla \vec{v})$ but this approximation can be useful when $J$ becomes inverted. \\

\textcolor{red}{To be added}: I forgot the proof that shows that adding the second term for approximation can help with inverted $J$. I will add it later.\\

\end{document}